\section{Test-Driven Development Strategy}

\subsection{Testing Philosophy}
The development process strictly follows Test-Driven Development (TDD) principles, where each module is developed only after its test suite is complete and passing. This approach ensures:

\begin{itemize}
    \item Complete test coverage before implementation
    \item Validation of both Pandas and DuckDB implementations
    \item Precise numerical accuracy in calculations
    \item Comprehensive edge case handling
\end{itemize}

\subsection{Test Data Engineering}

\subsubsection{Synthetic Dataset Design}
Each test suite uses carefully engineered test data:

\begin{lstlisting}[style=pythonstyle, caption=Test Data Example]
cls.test_data = pd.DataFrame({
    "date": ["2024-01-04", "2024-01-11"...],
    "user_id": [1, 1, 2, 2...],
    "total_calories": [2000, 2200...],
    "Aliment": ["Cacahuètes", "Amandes"...]
})
\end{lstlisting}

Key characteristics:
\begin{itemize}
    \item \textbf{Temporal Patterns:}
    \begin{itemize}
        \item Regular intervals (daily, weekly)
        \item Multiple weeks of data
        \item Specific weekday patterns
    \end{itemize}
    
    \item \textbf{User Behaviors:}
    \begin{itemize}
        \item Multiple users with different patterns
        \item Varied meal frequencies
        \item Different nutritional profiles
    \end{itemize}
    
    \item \textbf{Food Varieties:}
    \begin{itemize}
        \item Different food categories
        \item Various nutritional compositions
        \item Multiple portion sizes
    \end{itemize}
\end{itemize}

\subsection{Test Suite Architecture}

\subsubsection{Data Validation Tests}
\texttt{test\_1\_combine\_meal\_data.py}:

\begin{lstlisting}[style=pythonstyle, caption=Data Validation]
def test_calculated_values(load_dataframe):
    df_filtered = load_dataframe.dropna(subset=[
        "total_calories", "Valeur calorique"
    ])
    calories_match = np.isclose(
        df_filtered["total_calories"],
        df_filtered["Valeur calorique"] * 
        df_filtered["quantity"],
        rtol=1e-5
    )
\end{lstlisting}

Validation aspects:
\begin{itemize}
    \item Precise numerical calculations
    \item Proper handling of missing data
    \item Date/time format verification
    \item Data integrity checks
\end{itemize}

\subsubsection{Aggregation Tests}
\texttt{test\_3\_group\_data.py}:

\begin{itemize}
    \item \textbf{Test Data Structure:}
    \begin{itemize}
        \item Multiple days per user
        \item Various meal times
        \item Different food types
    \end{itemize}
    
    \item \textbf{Validation Points:}
    \begin{itemize}
        \item Daily totals accuracy
        \item User-specific aggregations
        \item Food type groupings
        \item Statistical computations
    \end{itemize}
\end{itemize}

\subsubsection{Window Function Tests}
Three specialized test modules:

\paragraph{Daily Windows} \texttt{test\_4\_1\_window\_function\_daily.py}:
\begin{itemize}
    \item Pre-calculated rolling averages
    \item 4-week patterns
    \item Weekday-specific calculations
\end{itemize}

\begin{lstlisting}[style=pythonstyle, caption=Window Function Test]
# Pre-calculated expected results
cls.expected_output["rolling_avg_total_calories"] = [
    2000.0,  # Initial value
    2100.0,  # Average of first two
    2080.0,  # Average of first three
    2150.0   # Average of all four
]
\end{lstlisting}

\paragraph{User Windows} \texttt{test\_4\_2\_window\_function\_user.py}:
\begin{itemize}
    \item User-specific patterns
    \item 4-day rolling windows
    \item Multiple metrics per user
\end{itemize}

\paragraph{Food Type Analysis} \texttt{test\_4\_3\_window\_function\_type\_food.py}:
\begin{itemize}
    \item Category proportions
    \item Nutritional distributions
    \item Time-based patterns
\end{itemize}

\subsection{Percentage Change Analysis}

\subsubsection{Daily Changes}
\texttt{test\_5\_1\_percentage\_change\_daily.py}:

\begin{lstlisting}[style=pythonstyle, caption=Percentage Change Test]
cls.expected_output["percentage_change_total_calories"] = [
    0.0,    # Base value
    4.76,   # (2200-2100)/2100 * 100
    -1.92,  # (2040-2080)/2080 * 100
    9.77    # (2360-2150)/2150 * 100
]
\end{lstlisting}

Test characteristics:
\begin{itemize}
    \item Manual calculation verification
    \item Edge case handling
    \item Precision control
\end{itemize}

\subsubsection{User Changes}
\texttt{test\_5\_2\_percentage\_change\_user.py}:
\begin{itemize}
    \item Individual baseline calculations
    \item Personal trend analysis
    \item Multi-user validation
\end{itemize}

\subsection{Testing Infrastructure}

\subsubsection{Base Classes}
Common testing infrastructure:
\begin{itemize}
    \item \texttt{TestDataSetup} for shared data
    \item \texttt{setUp} methods for each test class
    \item Cleanup procedures
    \item Error handling
\end{itemize}

\subsubsection{Assertion Methods}
Specialized validation:
\begin{itemize}
    \item \texttt{assert\_frame\_equal} with specific tolerances
    \item Date format verification
    \item Numerical precision checks
    \item Data integrity validation
\end{itemize}

\subsection{Development Process}

\subsubsection{TDD Workflow}
Step-by-step approach:
\begin{enumerate}
    \item Design test data with known outcomes
    \item Implement comprehensive test suite
    \item Verify test failures with empty implementation
    \item Develop feature to pass tests
    \item Refactor while maintaining test coverage
\end{enumerate}

\subsubsection{Quality Gates}
Implementation requirements:
\begin{itemize}
    \item All tests must pass
    \item Both implementations must match exactly
    \item Edge cases must be handled
    \item Performance metrics must be acceptable
\end{itemize}

This rigorous testing approach ensures the reliability and accuracy of the nutritional data pipeline, with particular emphasis on numerical precision, data integrity, and implementation consistency between Pandas and DuckDB approaches.
